% ======================= REFERENCIAS =======================
\clearpage
\section*{Referencias}
\begingroup
\small
\setlength{\parskip}{3pt}
\newcommand{\refe}[1]{\par\noindent\hangindent=1.5em\hangafter=1 #1\par}

\refe{Aguirre-López, M. A., Rodríguez-González, M. A., Soto-Villalobos, R., Gómez-Sánchez, L. E., Benavides-Ríos, Á. G., Benavides-Bravo, F. G., Walle-García, O., \& Pamanés-Aguilar, M. G. (2022). Statistical analysis of PM10 concentration in the Monterrey Metropolitan Area, Mexico (2010–2018). \emph{Atmosphere, 13}(2), Artículo 297. \url{https://doi.org/10.3390/atmos13020297}}

\refe{Bartlett, M. S. (1950). Tests of significance in factor analysis. \emph{British Journal of Statistical Psychology, 3}(2), 77–85. \url{https://doi.org/10.1111/j.2044-8317.1950.tb00285.x}}

\refe{Belachsen, I., \& Broday, D. M. (2022). Imputation of missing PM2.5 observations in a network of air quality monitoring stations by a new kNN method. \emph{Atmosphere, 13}(11), Artículo 1934. \url{https://doi.org/10.3390/atmos13111934}}

\refe{Box, G. E. P. (1949). A general distribution theory for a class of likelihood criteria. \emph{Biometrika, 36}(3/4), 317–346. \url{https://doi.org/10.1093/biomet/36.3-4.317}}

\refe{Brier, G. W. (1950). Verification of forecasts expressed in terms of probability. \emph{Monthly Weather Review, 78}(1), 1–3. \url{https://doi.org/10.1175/1520-0493(1950)078<0001:VOFEIT>2.0.CO;2}}

\refe{Carlstein, E. (1986). The use of subseries values for estimating the variance of a general statistic from a stationary sequence. \emph{The Annals of Statistics, 14}(3), 1171–1179. \url{https://doi.org/10.1214/aos/1176350057}}

\refe{Efron, B. (1975). The efficiency of logistic regression compared to normal discriminant analysis. \emph{Journal of the American Statistical Association, 70}(352), 892–898. \url{https://doi.org/10.1080/01621459.1975.10480319}}

\refe{Estévez, J., Gavilán, P., \& Giráldez, J. V. (2011). Guidelines on validation procedures for meteorological data from automatic weather stations. \emph{Journal of Hydrology, 402}(1–2), 144–154. \url{https://doi.org/10.1016/j.jhydrol.2011.02.031}}

\refe{Fawcett, T. (2006). An introduction to ROC analysis. \emph{Pattern Recognition Letters, 27}(8), 861–874. \url{https://doi.org/10.1016/j.patrec.2005.10.010}}

\refe{Fiebrich, C. A., Morgan, C. R., McCombs, A. G., Hall, P. K., \& McPherson, R. A. (2010). Quality assurance procedures for mesoscale meteorological data. \emph{Journal of Atmospheric and Oceanic Technology, 27}(10), 1565–1582. \url{https://doi.org/10.1175/2010JTECHA1433.1}}

\refe{Gobierno del Estado de Nuevo León. (2024, 7 de marzo). \emph{Busca Estado ganar atribuciones para sancionar a Refinería} [Boletín]. \url{https://nl.gob.mx/es/boletines/busca-estado-ganar-atribuciones-para-sancionar-refineria}}

\refe{Gómez, E., Sánchez, M., Medina, J., Montibeller, L., \& Pérez, S. (2025). \emph{MIMA-v1: Modelo de imputación para monitoreo ambiental} [Proyecto final del curso, no publicado]. Tecnológico de Monterrey.}

\refe{Gómez Santos, J. E., Sánchez Rodríguez, M., Medina García, J. E., Pérez Almazán, S., Montibeller González, J. L., Ruiz Hernández, B. R., \& Gómez Muñoz, R. (2027). A hybrid transformer-diffusion architecture for multivariate air quality data imputation. En H. Ponce, H. J. Escalante, \& R. A. Vazquez (Eds.), \emph{Artificial intelligence – COMIA 2026} (Communications in Computer and Information Science, Vol. 3097, pp. 728–739). Springer. \url{https://doi.org/10.1007/978-3-032-37441-7_55}}

\refe{Hadeed, S. J., O'Rourke, M. K., Burgess, J. L., Harris, R. B., \& Canales, R. A. (2020). Imputation methods for addressing missing data in short-term monitoring of air pollutants. \emph{Science of the Total Environment, 730}, Artículo 139140. \url{https://doi.org/10.1016/j.scitotenv.2020.139140}}

\refe{Hernández-Romero, K., Hernández-Romero, I. M., Mendoza, A., Pérez-Rodríguez, M., Barajas-Villarruel, L. R., Medellín-Carrillo, S., \& González, L. T. (2026). Persistent air pollution regimes and regulatory exceedances of PM10, PM2.5, and O3 in the Monterrey Metropolitan Area, Mexico. \emph{City and Environment Interactions, 31}, Artículo 100410. \url{https://doi.org/10.1016/j.cacint.2026.100410}}

\refe{James, G., Witten, D., Hastie, T., \& Tibshirani, R. (2021). \emph{An introduction to statistical learning: With applications in R} (2.ª ed.). Springer. \url{https://doi.org/10.1007/978-1-0716-1418-1}}

\refe{Johnson, R. A., \& Wichern, D. W. (2007). \emph{Applied multivariate statistical analysis} (6.ª ed.). Pearson Prentice Hall.}

\refe{Jolliffe, I. T., \& Cadima, J. (2016). Principal component analysis: A review and recent developments. \emph{Philosophical Transactions of the Royal Society A: Mathematical, Physical and Engineering Sciences, 374}(2065), Artículo 20150202. \url{https://doi.org/10.1098/rsta.2015.0202}}

\refe{Kaiser, H. F. (1960). The application of electronic computers to factor analysis. \emph{Educational and Psychological Measurement, 20}(1), 141–151. \url{https://doi.org/10.1177/001316446002000116}}

\refe{Kaiser, H. F., \& Rice, J. (1974). Little Jiffy, Mark IV. \emph{Educational and Psychological Measurement, 34}(1), 111–117. \url{https://doi.org/10.1177/001316447403400115}}

\refe{Krzanowski, W. J. (1975). Discrimination and classification using both binary and continuous variables. \emph{Journal of the American Statistical Association, 70}(352), 782–790. \url{https://doi.org/10.1080/01621459.1975.10480303}}

\refe{Künsch, H. R. (1989). The jackknife and the bootstrap for general stationary observations. \emph{The Annals of Statistics, 17}(3), 1217–1241. \url{https://doi.org/10.1214/aos/1176347265}}

\refe{Mardia, K. V. (1970). Measures of multivariate skewness and kurtosis with applications. \emph{Biometrika, 57}(3), 519–530. \url{https://doi.org/10.1093/biomet/57.3.519}}

\refe{Martínez-Cinco, M., Santos-Guzmán, J., \& Mejía-Velázquez, G. (2016). Source apportionment of PM2.5 for supporting control strategies in the Monterrey Metropolitan Area, Mexico. \emph{Journal of the Air \& Waste Management Association, 66}(6), 631–642. \url{https://doi.org/10.1080/10962247.2016.1159259}}

\refe{Mohd Shafi'i, M. S., \& Juahir, H. (2024). Assessment of spatial air quality on the East Coast of Peninsular Malaysia utilizing environmetric techniques. \emph{Environmental Monitoring and Assessment, 196}(7), Artículo 640. \url{https://doi.org/10.1007/s10661-024-12787-9}}

\refe{Moons, K. G. M., Donders, R. A. R. T., Stijnen, T., \& Harrell, F. E., Jr. (2006). Using the outcome for imputation of missing predictor values was preferred. \emph{Journal of Clinical Epidemiology, 59}(10), 1092–1101. \url{https://doi.org/10.1016/j.jclinepi.2006.01.009}}

\refe{Moritz, S., \& Bartz-Beielstein, T. (2017). imputeTS: Time series missing value imputation in R. \emph{The R Journal, 9}(1), 207–218. \url{https://doi.org/10.32614/RJ-2017-009}}

\refe{Orellano, P., Reynoso, J., Quaranta, N., Bardach, A., \& Ciapponi, A. (2020). Short-term exposure to particulate matter (PM10 and PM2.5), nitrogen dioxide (NO2), and ozone (O3) and all-cause and cause-specific mortality: Systematic review and meta-analysis. \emph{Environment International, 142}, Artículo 105876. \url{https://doi.org/10.1016/j.envint.2020.105876}}

\refe{Plaia, A., \& Bondì, A. L. (2006). Single imputation method of missing values in environmental pollution data sets. \emph{Atmospheric Environment, 40}(38), 7316–7330. \url{https://doi.org/10.1016/j.atmosenv.2006.06.040}}

\refe{Quinteros, M. E., Lu, S., Blazquez, C., Cárdenas-R, J. P., Ossa, X., Delgado-Saborit, J. M., Harrison, R. M., \& Ruiz-Rudolph, P. (2019). Use of data imputation tools to reconstruct incomplete air quality datasets: A case-study in Temuco, Chile. \emph{Atmospheric Environment, 200}, 40–49. \url{https://doi.org/10.1016/j.atmosenv.2018.11.053}}

\refe{Ramírez-Velarde, R., Esquivel-Flores, O., \& Mejía-Velázquez, G. (2024). Forecasting air pollution contingencies using predictive analytic techniques. \emph{Atmosphere, 15}(11), Artículo 1271. \url{https://doi.org/10.3390/atmos15111271}}

\refe{Rencher, A. C. (1992). Interpretation of canonical discriminant functions, canonical variates, and principal components. \emph{The American Statistician, 46}(3), 217–225. \url{https://doi.org/10.1080/00031305.1992.10475889}}

\refe{Roberts, D. R., Bahn, V., Ciuti, S., Boyce, M. S., Elith, J., Guillera-Arroita, G., Hauenstein, S., Lahoz-Monfort, J. J., Schröder, B., Thuiller, W., Warton, D. I., Wintle, B. A., Hartig, F., \& Dormann, C. F. (2017). Cross-validation strategies for data with temporal, spatial, hierarchical, or phylogenetic structure. \emph{Ecography, 40}(8), 913–929. \url{https://doi.org/10.1111/ecog.02881}}

\refe{Saminathan, S., \& Malathy, C. (2023). Ensemble-based classification approach for PM2.5 concentration forecasting using meteorological data. \emph{Frontiers in Big Data, 6}, Artículo 1175259. \url{https://doi.org/10.3389/fdata.2023.1175259}}

\refe{Schiavo, B., Morton-Bermea, O., Arredondo-Palacios, T. E., Meza-Figueroa, D., Robles-Morua, A., García-Martínez, R., Valera-Fernández, D., Inguaggiato, C., \& Gonzalez-Grijalva, B. (2023). Analysis of COVID-19 lockdown effects on urban air quality: A case study of Monterrey, Mexico. \emph{Sustainability, 15}(1), Artículo 642. \url{https://doi.org/10.3390/su15010642}}

\refe{Secretaría de Medio Ambiente del Estado de Nuevo León \& Clean Air Institute. (2024). \emph{Plan Integral de Gestión Estratégica de la Calidad del Aire 2023–2033 (PIGECA)}. Gobierno del Estado de Nuevo León. \url{https://aire.nl.gob.mx/docs/reportes/PIGECA_2023_2033.pdf}}

\refe{Secretaría de Salud. (2021, 27 de octubre). NOM-025-SSA1-2021, Salud ambiental. Criterio para evaluar la calidad del aire ambiente, con respecto a las partículas suspendidas PM10 y PM2.5. Valores normados para la concentración de partículas suspendidas PM10 y PM2.5 en el aire ambiente, como medida de protección a la salud de la población. \emph{Diario Oficial de la Federación}. \url{https://dof.gob.mx/nota_detalle_popup.php?codigo=5633855}}

\refe{Sistema Integral de Monitoreo Ambiental. (2026, 14 de agosto). \emph{Red de monitoreo y manejo de datos de calidad del aire} [Presentación de diapositivas]. Gobierno de Nuevo León.}

\refe{van Buuren, S., \& Groothuis-Oudshoorn, K. (2011). mice: Multivariate imputation by chained equations in R. \emph{Journal of Statistical Software, 45}(3), 1–67. \url{https://doi.org/10.18637/jss.v045.i03}}

\refe{Vélez-Pereira, A. M., Núñez-Magaña, N., Barreau, D., Bremer, K., \& O'Connor, D. J. (2025). Simplifying air quality forecasting: Logistic regression for predicting particulate matter in Chile. \emph{Atmosphere, 16}(12), Artículo 1377. \url{https://doi.org/10.3390/atmos16121377}}

\refe{Wilks, D. S. (2019). \emph{Statistical methods in the atmospheric sciences} (4.ª ed.). Elsevier.}

\refe{Wood, S. N. (2017). \emph{Generalized additive models: An introduction with R} (2.ª ed.). Chapman and Hall/CRC. \url{https://doi.org/10.1201/9781315370279}}

\refe{World Health Organization. (2021). \emph{WHO global air quality guidelines: Particulate matter (PM2.5 and PM10), ozone, nitrogen dioxide, sulfur dioxide and carbon monoxide}. \url{https://www.who.int/publications/i/item/9789240034228}}

\refe{Youden, W. J. (1950). Index for rating diagnostic tests. \emph{Cancer, 3}(1), 32–35. \url{https://doi.org/10.1002/1097-0142(1950)3:1<32::AID-CNCR2820030106>3.0.CO;2-3}}
\endgroup

% ======================= LIGA AL CÓDIGO =======================
\section*{Liga al código y organización}
\textbf{Repositorio:} \repo{} (\href{https://github.com/Leo-issacs/reto-sima-multivariados/tree/entrega-final-v2}{rama main, etiqueta \texttt{entrega-final-v2}}). Esta etiqueta agrega a \texttt{entrega-final} (el análisis entregado el 22 de octubre, que es el que enlaza la presentación) la verificación por temporada de la persistencia, la aplicación Shiny con el modelo final y esta versión del informe; el análisis principal no cambia. La etiqueta \texttt{informe-final} es la versión previa a la retroalimentación del 22 de octubre, y las decisiones de la etapa 4 se cerraron en \texttt{etapa4-decisiones-cerradas}. El README explica la organización y el orden de ejecución. La Tabla~\ref{tbl:codigo} relaciona los scripts con las secciones del informe.

\begin{table}[H]
\centering\footnotesize
\caption{Scripts del repositorio, en orden de ejecución, y sección del informe que alimentan.}
\label{tbl:codigo}
\begin{tabularx}{\textwidth}{>{\raggedright\arraybackslash}p{5.6cm}L>{\raggedright\arraybackslash}p{3.0cm}}
\toprule
Script & Qué hace & Sección\\
\midrule
\texttt{02\_importar\_diagnosticar.R} & Importa los 6 libros y diagnostica faltantes, rangos y consistencia & 4.1\\
\texttt{03\_limpiar.R}, \texttt{04\_validar\_limpios.R}, \texttt{05\_probar\_limpieza.R} & Limpia por reglas, agrega a diario y valida y prueba la limpieza & 4.2, Anexo B\\
\texttt{06\_explorar\_etapa2.R} & Construye y describe la muestra de la etapa 2 & 4.5\\
\texttt{07a}, \texttt{07b}, \texttt{07c} & Etapa 3: validación progresiva, evaluación en la prueba y diagnósticos & 5.7, 5.8, Anexo A\\
\texttt{08a\_enmascaramiento\_horario.R} & Elige la imputación horaria por enmascaramiento & 4.3\\
\texttt{08b\_construir\_base\_v2.R} & Construye la base v2 & 4.3, 4.5\\
\texttt{08c\_enmascaramiento\_diario.R}, \texttt{08d\_imputar\_diario.R} & Valida y aplica la imputación diaria (y MICE) & 4.3\\
\texttt{09a\_evaluar\_prueba\_v2.R} & Evaluación en la prueba con v2, tabla puente y sensibilidades & 5.7, 6, Anexo C\\
\texttt{09b\_validar\_lda.R} & Estabilidad de coeficientes, logística, pliegues, QDA y calibración & 5.4, 6.3, Anexo C\\
\texttt{09c\_area\_estudio\_reubicaciones.R} & Superación por estación, reubicaciones y mapa & 3\\
\texttt{09d\_diagnostico\_errores\_v2.R}, \texttt{09e\_sensibilidades\_posteriores.R} & Errores, supuestos y sensibilidades posteriores & 5.4, 5.6, 6.4, Anexo C\\
\texttt{09f\_viento\_pm\_no2.R} & Viento, PM2.5/PM10 y NO2 & 3\\
\texttt{09g\_pca\_y\_residuos\_v2.R} & PCA v2, residuos, largo de los huecos, radiación por estación y figuras compactas & 3, 5.2, 6.1, 6.2, Anexo B\\
\texttt{09h\_figuras\_informe.R} & Versión para el informe de siete figuras ya calculadas (solo presentación) & 3, 5.6, 6.2, Anexos B y C\\
\texttt{10\_analisis\_retro\_final.R} & Análisis posteriores a la retroalimentación final: logística, linealidad, persistencia, umbral, bootstrap estratificado, episodios y exclusiones & 4.5, 5.5, 5.7, 6.1, 6.3, 6.5, 6.6, Anexos B y C\\
\texttt{verificacion/}\allowbreak\texttt{persistencia\_por\_temporada.py} & Verificación: persistencia por temporada (Tabla~\ref{tbl:persistencia}) y correlaciones sin ajustar de la Sección~\ref{sec:coef} & 6.3, 6.5\\
\texttt{12\_exportar\_modelo\_app\_v2.R} & Exporta el modelo final para la aplicación Shiny (\texttt{app/}) & Liga al código\\
\bottomrule
\end{tabularx}
\end{table}

\textbf{Lenguajes.} R hace todo el análisis principal: limpieza, imputación, modelos, validaciones, figuras y la aplicación Shiny. Python se usó en la prueba de viabilidad del 4 de octubre, en verificaciones cruzadas y en el desglose por temporada de la persistencia (Tabla~\ref{tbl:persistencia}), que reproduce exactamente el total calculado en R. \textbf{Aplicación Shiny (\texttt{app/}).} Usa el modelo final (\texttt{output/etapa4/modelo\_final\_v2.rds}) y presenta la evaluación retrospectiva con meteorología observada; no es un pronóstico: para usarla como tal habría que evaluarla con meteorología pronosticada. \textbf{Cómo correrlo:} \texttt{renv::restore()} instala las versiones de los paquetes. Los scripts 02 y 03 necesitan los Excel de SIMA (en \texttt{data/raw/} o en la carpeta de \texttt{SIMA\_DATA\_DIR}); de 06 en adelante todo corre solo con \texttt{data/clean/}, que está en el repositorio. Se comprobó desde un clon limpio. El dataset de MIMA-v1 no se usa ni se redistribuye (licencia CC BY-ND); está en el anexo de Gómez et al. (2025).
